\subsection{Problem~1}%
\label{problem:1}

Find numbers $x_1$ and $x_2$ that maximize the sum $x_1 + x_2$ subject to the constraints:
\begin{equation*}
    \begin{split}
        x_1\geq 0, x_2\geq0,\\
        x_1+2x_2\leq4,\\
        4x_1+2x_2\leq12,\\
        -x_1+x_2\leq1,
    \end{split}
\end{equation*}

Draw the feasible set on $\Re^2$ and find the feasible solution geometrically and algebraically by
formulating a LP problem.

\subsubsection*{Mathematics}
Linear Programming (LP) is a type of optimization method used to minimize or maximize a linear objective function. A general LP problem has the following form:
\begin{equation*}
    \text{minimize}\quad c_1x_1+c_2x_2+...+c_nx_n=c^Tx
\end{equation*}
where $x \in \mathbb{R}^n$, $c \in \mathbb{R}^n$ which satisfy finitely many constrains of the form
\begin{equation*}
    \begin{split}
        a_{i1}x_1+a_{i2}x_2+...+a_{in}x_n\leq b_i \quad \text{(Inequality constraints)}\\
        a_{j1}x_1+a_{j2}x_2+...+a_{jn}x_n = b_j\quad \text{(Equality constraints)}\\
        a_{k1}x_1+a_{k2}x_2+...+a_{kn}x_n\geq b_k\quad \text{(Inequality constraints)}\\
    \end{split}
\end{equation*}
The function $c^Tx$ to be minimized is called \textit{objective function}. For each $x \in \mathbb{R}^n$ which satisfied all the constraints  is said to be a \textit{feasible point} for the problem. By introducing additional (slack) variables and equations, the linear programming problem can be put in a form in which the only constraints which appear are equalities or elementary inequalities. Which give us standard form :

\begin{equation*}
\begin{split}
    \min_{x}\ & \{c^Tx\} \\
    \text{subject to}\quad & Ax = b \quad \text{and} \quad x \geq 0 \\
    \text{where}\quad & x \in \mathbb{R}^n, \ c \in \mathbb{R}^n, \ A \in \mathbb{R}^{m \times n}
\end{split}
\end{equation*}

Where $\matr{A}^{m\times n}$ is a matrix of constrains and $\Vec{b}$ is a vector of solutions in $\mathbb{R}^m$.

\subsubsection*{The Two-Phase Simplex method}

The two-phase simplex method is a systematic approach for solving linear programming problems. It consists of two distinct phases that work sequentially to first establish feasibility and then find the optimal solution.
The goal is to find a feasible solution to the system.

\begin{equation*}
    \matr{A}x=\matr{b},\quad x_i \geq 0 \text{ for } i \in I.
\end{equation*}
\begin{enumerate}
    \item \textbf{Phase 1:}If no obvious feasible solution is available, an artificial variables are introduced to create an auxiliary problem.
    The objective of this auxiliary problem is to minimize the sum of these artificial variables.

    If the minimum value achieved is zero, it confirms that the original problem has at least one feasible solution, and the corresponding basis can be used as a starting point for Phase 2.

    \item \textbf{Phase 2:}In this step we proceeds to optimize the original objective function. The algorithm computes the basic solution and evaluates the reduced costs of non-basic variables.

    If all reduced costs satisfy the optimality conditions the current solution is optimal. Otherwise, the method selects an entering variable that can improve the objective function value. A direction of movement is computed, and if no feasible improvement direction exists, the problem is unbounded.

    When improvement is possible, a pivot operation exchanges a basic variable with the entering variable, updating the basis while maintaining feasibility. This process iterates until an optimal solution is found or unboundedness is detected.
\end{enumerate}

\subsubsection*{Solution}
We seek $z = \max(x_1+x_2)$.
Since all constraints are of the $\leq$ form and the RHS is non-negative, we introduce slack variables $x_3, x_4, x_5 \geq 0$ to convert the inequalities to equalities:

\begin{equation*}
    \left\{
    \begin{aligned}
        &x_1 \geq 0, x_2 \geq 0, x_3 \geq 0, x_4 \geq 0, x_5 \geq 0 \\
        &x_1 + x_2 + x_3 = 4 \\
        &4x_1 + 2x_2 + x_4 = 12 \\
        &-x_1 + x_2 + x_5 = 1
    \end{aligned}
    \right.
\end{equation*}
We construct initial tableau for simplex method:
\begin{equation*}
    \matr{f}=[-1 -1],\quad \matr{A}=\begin{bmatrix}
        1 & 1\\
        4 & 2\\
        -1 & 1\\
    \end{bmatrix},\quad \matr{b}=\begin{bmatrix}
        4\\12\\1
    \end{bmatrix}
\end{equation*}

\begin{table}[ht]
    \centering
    \begin{tabular}{c|c c|c c c|c}
         &$x_1$ & $x_2$ & $x_3$ & $x_4$ & $x_5$ & b \\
        \hline
        $s_1$ &1 & 1 & 1 & 0 & 0 & 4\\
        $s_2$ &4 & 2 & 0 & 1 & 0 & 12\\
        $s_3$ &-1 & 1 & 0 & 0 & 1 & 1\\
        \hline
        $z$ &-1 & -1 & 0 & 0 & 0 & 0\\
    \end{tabular}
\end{table}
Where $s_1, s_2, s_3$ are slack variables introduced to convert inequalities into equalities. The first step is to choose the column with the most negative objective-row entry ($x_1$) as the pivot column, then compute the minimum ratio test over rows with a positive coefficient in that column:
\begin{equation*}
    \begin{split}
    \matr{x_1}=\begin{bmatrix}
        1\\
        4\\
        -1\\
    \end{bmatrix},\quad s_1=\frac{4}{1}=4, \quad s_2= \frac{12}{4}=3\\
    \text{Pivot row}\rightarrow \matr{s_2}
    \end{split}
\end{equation*}
\begin{equation*}
    \begin{split}
    \begin{amatrix}{5}{1}
    1 & 1 & 1 & 0 & 0 & 4\\
    4 & 2 & 0 & 1 & 0 & 12\\
    -1 & 1 & 0 & 0 & 1 & 1\\
    \hline
    -1 & -1 & 0 & 0 & 0 & 0 \\
    \end{amatrix} \xrightarrow{\frac{s_2}{4}}
    \begin{amatrix}{5}{1}
    0 & 1.5 & 1.0 & -0.25 & 0 & 1\\
    1 & 0.5 & 0 & 0.25 & 0 & 3\\
    0 & 1.5 & 0 & 0.25 & 1 & 4\\
    \hline
    0 & -0.5 & 0 & 0.25 & 0 & 3 \\
    \end{amatrix}
    \end{split}
\end{equation*}
After finding minimal ratio we divide the pivot row and update other rows accordingly.
We repeat the same step for each column with negative coefficient left in the $z$(object row) row.
\begin{equation*}
\begin{split}
\begin{amatrix}{5}{1}
    0 & 1.5 & 1.0 & -0.25 & 0 & 1\\
    1 & 0.5 & 0 & 0.25 & 0 & 3\\
    0 & 1.5 & 0 & 0.25 & 1 & 4\\
    \hline
    0 & -0.5 & 0 & 0.25 & 0 & 3 \\
    \end{amatrix}\xrightarrow{\frac{s_1}{1.5}}\
    \begin{amatrix}{5}{1}
    0.0000 & 1.0000 & 0.6667 & -0.1667 & 0.0000 & 0.6667\\
    1.0000 & 0.0000 & -0.3333 & 0.3333 & 0.0000 & 2.6667\\
    0.0000 & 0.0000 & -1.0000 & 0.5000 & 1.000  & 3.0000\\
    \hline
    0.0000 & 0.0000 & 0.3333 & 0.1667 & 0 & 3.3333 \\
    \end{amatrix}
    \end{split}
\end{equation*}

All objective-row entries are non-negative, so the algorithm terminates.
The optimal solution read from the final tableau is:

\begin{equation*}
    \systeme{x_1=2.6667,x_2=0.6667}
\end{equation*}

\subsubsection*{Matlab Implementation}
\lstinputlisting[style=Matlab-editor]{problems/Problem_1.m}
